\paragraph{About this part.} The last part steps back from individual models. It asks what the
computing hardware allows, and why faster hardware cannot rescue a signal that decays too quickly
(Section~\ref{sec:hardware}); how models fail in live markets---regime changes, manipulation, hidden
liquidity and data errors (Section~\ref{sec:failures}); which open problems are most worth working on
(Section~\ref{sec:agenda}); and what the survey as a whole concludes (Section~\ref{sec:conclusion}).

% ================================================================
\section{Hardware acceleration and the infrastructure constraint}
\label{sec:hardware}
% ================================================================

Hardware interacts directly with both organising principles. The economic value of a signal
depends on the latency and the determinism of the infrastructure that evaluates it, and richer
models tend to raise both latency and jitter:
\begin{equation}
  \text{model complexity} \uparrow \;\Rightarrow\; \text{inference latency} \uparrow \;\Rightarrow\;
  \text{signal value at execution} \downarrow .
\end{equation}
\inwords{a bigger model takes longer to evaluate, and by the time it answers, some of what it
predicted has already happened.}

\subsection{Determinism as a risk-control property}

General-purpose processors, even with kernel-bypass networking, introduce jitter from caches,
branch prediction, interrupts and scheduling. A strategy that is profitable at its average latency
can lose money if latency spikes during the volatile episodes when its signal matters most.
Deterministic, fixed-cycle execution is therefore a risk-control property, not only a speed
optimisation.

\subsection{Book-building latency versus inference latency}

Two pipelines are often conflated. \emph{Book building}---parsing market data and maintaining the
order book---is protocol-bound and highly regular. \emph{Inference}---evaluating a model on the
book---depends on the model. Field-programmable gate arrays have long been used for the first:
feed handling and book maintenance in hardware are well documented~\citep{leber2011,lockwood2012}.
Placing models on the same fabric is more recent. Tool flows for fast neural-network inference on
FPGAs~\citep{hls4ml} make small, quantised networks feasible with very low and fixed latency,
but model capacity is bounded by the fabric. Vendor-audited benchmarks of FPGA trigger paths now
report minimum trigger latencies on the order of ten nanoseconds~\citep{exegy_stac2024}, and
hardware-aware sequence models designed for FPGA execution from the outset, rather than ported after
training, have been proposed for cryptocurrency tick data~\citep{popov2026ssm}. The two numbers
measure different pipelines---a trigger versus a learned model---and should not be compared
directly.

\subsection{The commercial FPGA landscape}

Most production hardware for trading is commercial, and most of its published figures are vendor
claims rather than audited measurements. Table~\ref{tab:vendors} lists representative offerings with
the figures their makers publish; the one independently audited result is the STAC-T0 benchmark of
an Exegy framework on an AMD card~\citep{exegy_stac2024}. Two cautions apply when reading such
numbers. They measure different things---a transceiver loopback, a trigger, a book build, a model
inference---that must not be compared directly. And the AMD transceiver figure explicitly excludes
protocol processing and programmable logic, so it is a lower bound on a component, not a
tick-to-trade time.

\begin{table}[ht]
\centering\footnotesize
\caption{Representative commercial hardware for trading and the figures published by the vendors.
Except for the STAC result, figures are vendor claims without independent audit.}
\label{tab:vendors}
\begin{tabularx}{\linewidth}{>{\raggedright\arraybackslash}p{0.2\linewidth}>{\raggedright\arraybackslash}X>{\raggedright\arraybackslash}X}
\toprule
Offering & What it is & Published figure \\
\midrule
AMD Alveo UL3524 / UL3422~\citep{amd_ul3524,amd_ul3422} & FPGA cards built for trading, 64 low-latency transceivers & under 3\,ns transceiver latency, raw loopback excluding logic and protocol \\
Exegy (incl.\ Enyx, acquired 2022)~\citep{exegy_enyx,exegy_stac2024} & FPGA feed handling, order entry and trigger frameworks & STAC-T0 audited trigger latency of the order of ten nanoseconds \\
Algo-Logic~\citep{algologic} & CME tick-to-trade system on Alveo U50, with feed handler, book and TCP order entry & under 500\,ns tick-to-trade \\
Magmio~\citep{magmio} & C++ strategies compiled to FPGA by high-level synthesis & under 100\,ns for triggers; under 350\,ns with book building \\
Xelera Silva~\citep{xelera_silva} & tree-ensemble and small-network inference on FPGA & about 250\,ns (99th percentile) for a gradient-boosted model inline on UL3524; 1.1--1.3\,\textmu s as card offload \\
Napatech + Xelera~\citep{napatech} & the same inference on FPGA network cards & 1.55\,\textmu s for a 1000-tree model with 100 features \\
\bottomrule
\end{tabularx}
\end{table}

\subsection{Between FPGA and GPU}

A class of accelerators aims at the gap between FPGA determinism and GPU flexibility.
LightTrader~\citep{lighttrader}, from KAIST and the chip designer Rebellions, combines an FPGA for
market data and orders with dedicated AI processors and a scheduler that manages workload and power;
back-tested on CME data, it reports speed-ups of about 14 times over a GPU-based and 7 times over an
FPGA-based system for AI inference. Corsair~\citep{corsair2025} is an in-memory-computing chiplet
design for large-language-model inference; it is not a trading product, but it illustrates the
direction of the hardware industry. Network cards with programmable logic carry inference into the
network path itself~\citep{napatech}.

\subsection{Transformers on FPGAs}

Mapping attention to FPGAs is an active research topic outside finance. FAMOUS~\citep{famous}
accelerates dense multi-head attention on a data-centre FPGA and reports about 2.6 times the throughput
of a V100 GPU; ELiTeFormer~\citep{eliteformer} co-designs a ternary-weight Transformer with linear
attention for an adaptive SoC. Neither is a trading system. For LOB models the implication is that
small, quantised attention models can be placed on programmable logic, at latencies of microseconds
rather than nanoseconds, while full spatial-temporal Transformers remain GPU workloads.
Quantisation---integer rather than floating-point arithmetic---is the enabling
technique~\citep{jacob2018}.

\subsection{GPUs: throughput for research}

Graphics processors excel at data-parallel workloads: training, large-scale simulation, and models
whose memory footprint or operator mix exceeds practical FPGA budgets. JAX-LOB~\citep{frey2023}
moves LOB matching onto accelerators for reinforcement-learning training. Monte Carlo simulation
of order-book dynamics parallelises well across paths; one vendor demonstration reports speed-ups of
14 to 114 times over a CPU~\citep{nvidia_numba2025}, though for simulated paths rather than a
historical matching-engine backtest. Specialised processors also accelerate training: a multi-horizon
DeepLOB variant trained in about 15\% of the GPU time on intelligence processing
units~\citep{zhangzohren2021}. For single-event,
latency-critical inference, batch-oriented accelerators are at a disadvantage; for signals whose
useful life is measured in milliseconds or longer they are adequate.

Faster research hardware has the effect noted in Section~\ref{sec:search}: it increases the number
of hypotheses that can be tested, and so the certification threshold must scale with it.

\subsection{Networks and distance}

Latency is not only computation. The operating system's network stack adds microseconds;
\emph{kernel-bypass} libraries deliver packets directly to user space~\citep{dpdk}. Published figures for
such stacks include half-round-trip times of about 3\,\textmu s on older 10-gigabit
hardware~\citep{openonload} and about 1\,\textmu s for UDP~\citep{vma}. Physical distance sets a floor no
software can remove: light in optical fibre takes about 5\,\textmu s per kilometre (about 4.9\,ns per
metre)~\citep{corning_smf28,fiberlatency}, which is why trading systems are colocated in the
exchange's data centre and why differences in relative speed translate into differences in trading
performance~\citep{baron2019,budish2015}. At the extreme, application-specific chips fix a design in
silicon, trading the reprogrammability of an FPGA for speed and power~\citep{kuon2007}.

\subsection{Workload placement}

\begin{table}[ht]
\centering\small
\caption{A qualitative guide to workload placement. It reflects the general division of labour
between platforms, not measured results.}
\label{tab:placement}
\begin{tabularx}{\linewidth}{lXX}
\toprule
Workload & Preferred platform & Rationale \\
\midrule
Market-data parsing, book maintenance & FPGA & regular, protocol-bound, deterministic \\
Linear OFI, queue imbalance, simple triggers & FPGA or CPU & few operations; easy to pipeline \\
Small quantised networks or tree ensembles on the critical path & FPGA & feasible when capacity is
  constrained by design \\
Pre-trade risk checks, mass cancel & FPGA & must be deterministic and next to order generation \\
Multi-horizon or cross-sectional deep models, full Transformers & GPU or other accelerator &
  memory and operator mix exceed FPGA budgets \\
Large-scale simulation, training, hyper-parameter search & GPU & throughput dominates \\
Strategy logic that changes often & CPU with kernel bypass & flexibility outweighs nanoseconds \\
\bottomrule
\end{tabularx}
\end{table}

The guiding rule is Occam's razor (Principle~3) again: place a computation on the most constrained, and
therefore most deterministic, hardware that still delivers the required economic performance.
Moving a model to faster hardware is justified only when the reduction in latency and jitter
measurably improves post-cost performance.

\subsection{Hardware cannot extend signal lifetime}

Simpler models have a structural advantage on deterministic hardware: fewer parameters and
simpler data flow map onto fixed-point, branch-free pipelines. Acceleration can rescue part of a
complex model's value by removing engineering latency, but it cannot extend the statistical
lifetime of the signal.

\begin{example}[Second-order latency]
A logistic score on five OFI windows and queue imbalance takes a few tens of multiply--add
operations. If the signal's half-life is 5 ms, evaluating it in 50\,$\mu$s rather than 500\,ns
keeps $2^{-0.01} \approx 99.3\%$ of its value rather than essentially all of it. The difference is
second-order; the binding constraint is the quality and horizon of the signal itself.
\end{example}

The design sequence that follows is: identify the horizon and value of the signal under idealised
execution; choose the simplest model that captures that value; only then choose the hardware tier
whose latency and jitter preserve it after costs.

\subsection{Open questions in hardware}

\begin{itemize}[leftmargin=1.5em,itemsep=1pt]
\item Training objectives that penalise expected inference latency and resource use on a target
  device, so that models trade accuracy against deployability explicitly.
\item Which components of spatial-temporal LOB models survive mapping to programmable logic without
  losing their economic value.
\item Hawkes intensities in hardware: exponential kernels update in constant time and suit pipelines,
  and a power-law kernel can be approximated by a few exponentials, each with its own running
  sum~\citep{bochud2007}; how many are needed at the precision of fixed-point hardware, and how
  state-dependent kernels map to pipelines, are open.
\item How certification thresholds should scale when faster simulation multiplies the number of
  candidates that can be tested.
\item Whether the statistical properties of a strategy developed on one hardware and latency profile
  survive deployment on another---the research--production gap of Part~IV, at the hardware layer.
\item Public, auditable benchmarks that cover book building and inference together, rather than
  isolated kernels or vendor-defined paths.
\end{itemize}

% ================================================================
\section{Failure modes}
\label{sec:failures}
% ================================================================

Models are fitted to the past and run in the future. This section lists the ways that gap breaks
them in practice: the market changes regime, liquidity disappears in stress, other participants place
orders designed to mislead, displayed size hides real size, and the data themselves are wrong. For each
it states what goes wrong and, where one exists, how it can be detected or guarded against.

\paragraph{Regime change and state invalidation.} A model fitted in one liquidity regime can
fail after volatility shocks, tick-size or fee changes, market-structure changes, scheduled
announcements or liquidity withdrawal. A signal can be statistically valid for the regime in which
it was estimated and economically stale after a transition.

\paragraph{Volatility and liquidity withdrawal.} Queue-based signals rely on the book looking roughly
the same when a decision is executed as when the signal was computed. In a volatility burst, best-level
depth can vanish and the intensities that drive depletion change on the same time scale as the queues
themselves, so first-passage estimates become unreliable. A simple defence is to reduce size or disable
microstructure signals when short-term realised volatility is unusually high; it should be compared
against more elaborate regime models before those are adopted.

\paragraph{Spoofing and layering.} Displayed liquidity is not necessarily committed liquidity.
Manipulative orders do not falsify OFI---OFI still measures observed changes---but they alter the
mapping from displayed flow to future executable prices, and they distort queue-imbalance, depth
and depletion estimates.

\paragraph{Documented cases.} Manipulation of the order book is not hypothetical. On the Korea Exchange,
\citet{lee2013spoof} document traders placing large orders away from the market to create imbalance and
withdrawing them. A trader in the E-mini S\&P~500 was charged with layering over five years, conduct the
CFTC alleged contributed to the conditions of the May 2010 flash crash~\citep{cftc2015sarao}. In 2020
JPMorgan paid \$920 million over spoofing in precious-metals and US Treasury futures over at least eight
years~\citep{cftc2020jpm}. In a spoofing episode in Ultra Treasury bond futures, \citet{debie2023} measure ask-side
liquidity costs falling from about 10.4 to about 8 basis points when the spoof order was placed and
recovering after it was cancelled---a displayed improvement in liquidity that was not real.
\citet{cartea2020spoof} model spoofing as an optimal-control problem in which the spoofer trades the gain
from skewing the order imbalance against the expected regulatory penalty. Quote stuffing---bursts of
order messages---affected over 74\% of US-listed securities at least once in 2010, with lower liquidity
and higher short-term volatility during episodes~\citep{egginton2016}.

\paragraph{Hidden liquidity.} Iceberg orders make executable liquidity exceed displayed liquidity,
so a model trained on displayed volume misestimates depletion, impact and fill probability. In
E-mini S\&P~500 futures, one detection study estimates that about 9\% of book volume is iceberg
orders~\citep{christensen2013}; when participants detect an iceberg they trade against it actively, and
icebergs create an adverse-selection cost for other limit orders~\citep{frey2017iceberg}. This is one
reason message-level and trade data are valuable even for book-based models.

\begin{table}[ht]
\centering\footnotesize
\caption{How contamination affects order-book signals.}
\label{tab:contamination}
\begin{tabularx}{\linewidth}{l>{\raggedright\arraybackslash}X>{\raggedright\arraybackslash}X}
\toprule
Form & Signature in the data & Effect on signals \\
\midrule
Spoofing & large orders away from the touch, cancelled quickly & OFI and imbalance point the wrong way \\
Layering & several non-genuine orders at successive levels & distorts depth, slope and multi-level OFI \\
Quote stuffing & bursts of submissions and cancellations & inflates event counts; breaks event-clock assumptions \\
Icebergs and hidden orders & repeated executions at one price with stable visible size & displayed depth understates liquidity; depletion and fill estimates biased \\
\bottomrule
\end{tabularx}
\end{table}

Mitigations include down-weighting orders that are cancelled within a short time, features built from
cancellation rates and from the ratio of trades to orders, and detection of iceberg refills. They reduce
but do not remove the problem; robustness to each form of contamination belongs in a serious
evaluation.

\paragraph{Feed and timestamp errors.} At high frequency, data quality is part of the model:
timestamp synchronisation, packet loss, sequence gaps, book reconstruction, exchange-specific event
semantics, stale snapshots, duplicates, and clock differences between market-data and execution
systems.

\begin{example}[A sequence gap]
On a feed with sequence numbers, a missed packet leaves the reconstructed book wrong until the
handler recovers from a snapshot. A model trained on rows from that interval learns
reconstruction error, not market behaviour. No number of parameters compensates for it; the gap
has to be detected and the interval excluded or repaired.
\end{example}

% ================================================================
\section{Research agenda}
\label{sec:agenda}
% ================================================================

The survey has identified several places where the literature stops short of what a trader needs.
This section collects them as open research questions; Part~VI gives the standard against which
answers to them should be evaluated.

\paragraph{Execution-aware objectives.} Most prediction models stop at
an estimate of $\P(\Delta m_{t,H} > 0 \mid X_t)$. The more useful objective is an estimate of
$\E[\Pi \mid a_t, X_t]$, the expected profit of each available action,
combining prediction, fill probability, queue position and adverse selection. Calibrated
confidence~\citep{manoharan2026uqlob} is a step toward it: it tells an execution policy when to act.

\paragraph{Queue-aware benchmarks.} Queue position should be a first-class state variable.
Benchmarks should report performance conditional on queue rank, residual volume, spread and expected
depletion time, and should measure the size of the queue-rank effect on fill probability and on
conditional mark-outs across assets and regimes, rather than take a single model estimate as
given.

\paragraph{Message-level foundation models for execution.} Pretrained message-level encoders have
so far been evaluated mainly on mid-price and next-message prediction~\citep{linna2025lobert}.
Fine-tuning them for fill probability and conditional mark-outs, evaluated in queue-exact replay,
is an open direction.

\paragraph{Cross-asset information after costs.} When does cross-asset information remain valuable
after costs, latency, multiple testing and parameter instability?

\paragraph{Robustness to structural change.} Models should be tested across volatility regimes,
tick and fee changes, market-structure changes and extreme events.

\paragraph{Automated search as a statistical process.} Automated and LLM-assisted search should be
evaluated as a search: with complete trial accounting, search-aware nulls, and a reported
false-discovery rate. A public, LOB-specific skill-bar benchmark---data with null features and with
signals of known strength injected---would let competing search systems be compared on how often they
find what is not there and how often they miss what is.

\paragraph{End-to-end signal and execution.} Most models output a forecast and leave execution to a
separate rule. Models whose output is the decision itself---quote, keep, cancel, cross---and whose
training objective is execution value would close the gap that this survey describes. Reinforcement
learning in LOB simulators~\citep{lalor2025nhp,frey2023} is one route, ideally validated in queue-exact replay; its results must be
judged against model-free benchmarks in the same simulator.

\paragraph{Causal rather than predictive questions.} The questions that matter for execution are
causal: what happens to the queue and the price \emph{if} an order is placed here? Replay with a
no-impact assumption answers this only for small orders. Models of the market's response to one's own
orders, and ways to validate them without live experiments, remain open.

\paragraph{State-space and long-context models for event streams.} State-space sequence models are well
suited to long, irregular event streams, yet apart from one hardware-oriented design~\citep{popov2026ssm} we found no established LOB forecaster built on them.
Testing them against the baselines of Table~\ref{tab:selection}---not against each other---is a clear
next step.

\paragraph{Near-criticality and learned models.} Estimated branching ratios of order flow are high,
from about $0.7$--$0.8$ to about one depending on method. How deep sequence
models behave on near-critical, long-memory inputs, and whether Hawkes structure built into them
(Section~\ref{sec:hybrid}) improves robustness, is untested.

\paragraph{Contamination-robust features.} Features that remain informative in the presence of spoofing
and hidden liquidity---based on order lifetimes, cancellation behaviour and refill patterns---and
evaluation protocols that test robustness explicitly.

\paragraph{Generative evaluation.} Generators should be evaluated by market invariants and downstream
utility rather than by visual or marginal similarity alone~\citep{lobbench2025}.

\paragraph{A parsimony audit.} For each widely cited deep LOB model, re-estimate its incremental
economic value and deflated performance against the best low-capacity baseline on the same data and
fill model. Such a meta-study would directly test the central claim of this survey.

\paragraph{Correlated signals and systemic risk.} If many participants deploy similar order-book signals,
their trades become correlated, and so do their withdrawals of liquidity under stress. The market-wide
consequences of widely shared LOB alphas are largely unstudied.

% ================================================================
\section{Conclusion}
\label{sec:conclusion}
% ================================================================

The literature on limit-order-book modelling has moved from low-dimensional queue and order-flow
models to expressive sequence architectures and pretrained message-level encoders. Statistical
prediction has improved. Statistical prediction and executable trading remain different problems.
The translation between them is
\[
  \boxed{\ \text{LOB state} \to \text{representation} \to \text{prediction} \to \text{decision} \to
  \text{queue} \to \text{execution} \to \text{P\&L}\ }
\]
and each link can destroy information or add uncertainty.

Order-flow imbalance, queue imbalance, Hawkes intensities and first-passage models give
interpretable low-dimensional descriptions of market dynamics. Deep and pretrained models give
representational capacity and, increasingly, calibrated confidence. Hybrid models offer a route to
combine them. The conclusion is not that simple models are superior. It is that complexity is an
empirical hypothesis: a more complex model is justified when it delivers incremental out-of-sample
economic value that survives realistic execution, temporal validation, and the multiple-testing
burden of its own selection.

Two tools make that standard testable. A deterministic, queue-exact replay turns the execution
questions into measurements on observed data. A disciplined model-discovery loop, in which a
language model proposes and deterministic machinery disposes, widens the search without weakening
the evidence---provided the trial count is part of the result. Automation raises the rate at which
hypotheses are generated. It does not remove the statistical problem of deciding whether a
discovered pattern is real.

For a reader who must choose a model, the survey's guidance reduces to an order of operations.
Decide first what the output is for---aggressive trading, passive execution, market making or
simulation---and therefore which of the targets of Section~\ref{sec:deep} matters. Establish the
strongest simple baseline on the best available representation: order-flow imbalance, queue
imbalance, the micro-price, the queue race, and for passive orders the volume ahead. Evaluate it with
proper scores and economic metrics in a queue-exact replay, counting every configuration tried. Only
then add capacity---a small network, a sequence model, a hybrid---and keep it only if the increment
survives the same test (Table~\ref{tab:selection}).

A useful future literature should therefore report not only whether a model predicts the market,
but what information it extracts, what decision it changes, how that decision is executed, and
whether the incremental value survives a properly constructed null.

\begin{table}[ht]
\centering\small
\caption{Evidence status of the main claims in this survey.}
\label{tab:evidence}
\begin{tabularx}{\linewidth}{Xl}
\toprule
Claim & Status \\
\midrule
Short-horizon mid changes are approximately linear in OFI, with slope inversely related to
  depth~\citep{cks2014} & established \\
Hawkes stationarity when $\rho(\Gamma) < 1$ and $\bar\lambda = (I-\Gamma)^{-1}\mu$ & established \\
Input representation can matter as much as the architecture & synthesis \\
Naive ``touch'' backtests overstate passive fills relative to queue-exact replay & synthesis;
  argued from price--time priority \\
Front-of-queue positions are worth more than back-of-queue positions, with adverse selection
  increasing in queue position~\citep{moallemi2017} & established in a model; size across regimes open \\
Hybrid inputs derived from the same stream add little to a sufficiently large network & open
  hypothesis \\
LLM-driven search raises the multiple-testing burden & established in principle \\
\bottomrule
\end{tabularx}
\end{table}
